% GENERATED FILE. Edit canonical lessons, then run npm run generate:books.
% canonical-source-hash: fnv1a64:a4c3f246056f650c
% canonical-lessons: SA-C127-daridrah, SA-C127-dhanikah, SA-C127-tarunah, SA-C127-prasiddhah, SA-C127-suraksitah

\chapter{Poor, Rich, Young, Famous, Safe}
\label{ch:sa-poor-rich-young-famous-safe}
\hlchaptermodality{127}

\begin{chapteropening}
\textbf{By the end of this chapter:} \emph{I can say poor, rich, young, famous and safe.}

Complete the last lesson of chapter 127: i can say poor, rich, young, famous and safe.
\end{chapteropening}

\section[\texorpdfstring{\sk{दरिद्रः} (daridraḥ)}{दरिद्रः (daridraḥ)}]{\sk{दरिद्रः} (daridraḥ) — poor}

\label{lesson:SA-C127-daridrah}



\begin{quote}
Before the new one: say the Sanskrit for wet, then the Sanskrit for dry.
\end{quote}

\subsection*{You'll want to know: \sk{दरिद्रः}}
\textbf{\sk{दरिद्रः}} — \emph{daridraḥ} — \textquotedblleft{}poor\textquotedblright{}.

\begin{grammarlens}[title={What you've built}]
One more word for this chapter.
\end{grammarlens}

\subsection*{Your turn}
\begin{itemize}
\raggedright
  \item \emph{Say it:} \emph{daridraḥ}
  \item \emph{Say it:} \emph{daridraḥ}, once more
  \item \emph{Say it:} say \emph{daridraḥ}
  \item \emph{Recall:} say the Sanskrit for wet, then the Sanskrit for dry, then say \emph{daridraḥ} again
\end{itemize}

\begin{tcolorbox}[breakable,colback=teal!4,colframe=teal!35!black,title={Before you move on}]
What does \sk{दरिद्रः} mean? (\textquotedblleft{}Poor\textquotedblright{}.) Say it once more.
\end{tcolorbox}
\section[\texorpdfstring{\sk{धनिकः} (dhanikaḥ)}{धनिकः (dhanikaḥ)}]{\sk{धनिकः} (dhanikaḥ) — rich}

\label{lesson:SA-C127-dhanikah}



\begin{quote}
Before the new one: say the Sanskrit for dry, then the Sanskrit for poor.
\end{quote}

\subsection*{You'll want to know: \sk{धनिकः}}
\textbf{\sk{धनिकः}} — \emph{dhanikaḥ} — \textquotedblleft{}rich\textquotedblright{}.

\begin{grammarlens}[title={What you've built}]
One more word for this chapter.
\end{grammarlens}

\subsection*{Your turn}
\begin{itemize}
\raggedright
  \item \emph{Say it:} \emph{dhanikaḥ}
  \item \emph{Say it:} \emph{dhanikaḥ}, once more
  \item \emph{Say it:} say \emph{dhanikaḥ}
  \item \emph{Recall:} say the Sanskrit for dry, then the Sanskrit for poor, then say \emph{dhanikaḥ} again
\end{itemize}

\begin{tcolorbox}[breakable,colback=teal!4,colframe=teal!35!black,title={Before you move on}]
What does \sk{धनिकः} mean? (\textquotedblleft{}Rich\textquotedblright{}.) Say it once more.
\end{tcolorbox}
\section[\texorpdfstring{\sk{तरुणः} (taruṇaḥ)}{तरुणः (taruṇaḥ)}]{\sk{तरुणः} (taruṇaḥ) — young}

\label{lesson:SA-C127-tarunah}



\begin{quote}
Before the new one: say the Sanskrit for poor, then the Sanskrit for rich.
\end{quote}

\subsection*{You'll want to know: \sk{तरुणः}}
\textbf{\sk{तरुणः}} — \emph{taruṇaḥ} — \textquotedblleft{}young\textquotedblright{}.

\begin{grammarlens}[title={What you've built}]
One more word for this chapter.
\end{grammarlens}

\subsection*{Your turn}
\begin{itemize}
\raggedright
  \item \emph{Say it:} \emph{taruṇaḥ}
  \item \emph{Say it:} \emph{taruṇaḥ}, once more
  \item \emph{Say it:} say \emph{taruṇaḥ}
  \item \emph{Recall:} say the Sanskrit for poor, then the Sanskrit for rich, then say \emph{taruṇaḥ} again
\end{itemize}

\begin{tcolorbox}[breakable,colback=teal!4,colframe=teal!35!black,title={Before you move on}]
What does \sk{तरुणः} mean? (\textquotedblleft{}Young\textquotedblright{}.) Say it once more.
\end{tcolorbox}
\section[\texorpdfstring{\sk{प्रसिद्धः} (prasiddhaḥ)}{प्रसिद्धः (prasiddhaḥ)}]{\sk{प्रसिद्धः} (prasiddhaḥ) — famous}

\label{lesson:SA-C127-prasiddhah}



\begin{quote}
Before the new one: say the Sanskrit for rich, then the Sanskrit for young.
\end{quote}

\subsection*{You'll want to know: \sk{प्रसिद्धः}}
\textbf{\sk{प्रसिद्धः}} — \emph{prasiddhaḥ} — \textquotedblleft{}famous\textquotedblright{}.

\begin{grammarlens}[title={What you've built}]
One more word for this chapter.
\end{grammarlens}

\subsection*{Your turn}
\begin{itemize}
\raggedright
  \item \emph{Say it:} \emph{prasiddhaḥ}
  \item \emph{Say it:} \emph{prasiddhaḥ}, once more
  \item \emph{Say it:} say \emph{prasiddhaḥ}
  \item \emph{Recall:} say the Sanskrit for rich, then the Sanskrit for young, then say \emph{prasiddhaḥ} again
\end{itemize}

\begin{tcolorbox}[breakable,colback=teal!4,colframe=teal!35!black,title={Before you move on}]
What does \sk{प्रसिद्धः} mean? (\textquotedblleft{}Famous\textquotedblright{}.) Say it once more.
\end{tcolorbox}
\section[\texorpdfstring{\sk{सुरक्षितः} (surakṣitaḥ)}{सुरक्षितः (surakṣitaḥ)}]{\sk{सुरक्षितः} (surakṣitaḥ) — safe}

\label{lesson:SA-C127-suraksitah}



\begin{quote}
Before the new one: say the Sanskrit for young, then the Sanskrit for famous.
\end{quote}

\subsection*{You'll want to know: \sk{सुरक्षितः}}
\textbf{\sk{सुरक्षितः}} — \emph{surakṣitaḥ} — \textquotedblleft{}safe\textquotedblright{}.

\begin{grammarlens}[title={What you've built}]
One more word for this chapter.
\end{grammarlens}

\subsection*{Your turn}
\begin{itemize}
\raggedright
  \item \emph{Say it:} \emph{surakṣitaḥ}
  \item \emph{Say it:} \emph{surakṣitaḥ}, once more
  \item \emph{Say it:} say \emph{surakṣitaḥ}
  \item \emph{Recall:} say the Sanskrit for young, then the Sanskrit for famous, then say \emph{surakṣitaḥ} again
\end{itemize}

\begin{tcolorbox}[breakable,colback=teal!4,colframe=teal!35!black,title={Before you move on}]
What does \sk{सुरक्षितः} mean? (\textquotedblleft{}Safe\textquotedblright{}.) Say it once more.
\end{tcolorbox}
